\scanpage{58}
% archive sheet number: 57
\deleted{\uline{Proposition}} \inserted{Lemme} $S$ \deleted{\unsure{schéma en groupes noeth}} \inserted{préschém. loc. noeth},
\margin{\rule[-2\baselineskip]{0.35pt}{2\baselineskip}}$G$, \inserted{$G_1$,} préschém. en groupes \inserted{de type fini} sur $S$, \deleted{\unsure{et}}
plats sur $S$, $u:G\unsure{\longrightarrow|\text{---}}G_1$ un morphisme \inserted{\unsure{de groupe|de groupes|à fibres}}
\uline{surjectif}, (donc nécessairement plat, \unsure{réf|cf.|surj.}),
% slanted
\margin{conditions de finitude inutile, si on suppose à priori $u$ plat \unsure{---}}$N$ son noyau (qui est \inserted{donc} plat sur $S$).
Alors il y a correspondance biunivoque
\margin{\rule[-10\baselineskip]{0.35pt}{10\baselineskip}}entre les sous-schémas fermés \inserted{$H$} de
$G$ qui sont invariants par $\deleted{\unsure{H}}N$, et les
sous schémas fermés \inserted{$H_1$} de \deleted{\unsure{$k$}} $G_1$, à $H_1$
correspondant $H=u^{-1}(H_1)$. Pour que
$H$ soit \deleted{propre} plat sur $S$, il f. et s. que
$H_1$ le soit. Pour que $H$ soit un
ss-groupe de $G$, il faut et il suffit que
$H_1$ soit un ss-groupe de $G_1$. Pour
que \inserted{de plus} $H_1\longrightarrow$ soit radiciel surjectif
sur $S$, il f. et s. que $H$ et $N$ aient
$=$ \deleted{\unsure{esp.}} \inserted{ens.} sous-jacent. \deleted{Alors} (Comme
$N\longrightarrow S$ \inserted{(donc $H\longrightarrow S$)} est \inserted{universelt} ouvert) c'est un morphisme
\uline{fini} radiciel \uline{surjectif}.

AQT par \uuline{théorie de descente}.
